\documentclass[12pt,a4paper]{article}
\usepackage[a4paper,margin=21mm]{geometry}
\usepackage[T2A]{fontenc}
\usepackage[utf8]{inputenc}
\usepackage[russian]{babel}
\usepackage{amsmath,amssymb,mathtools}
\usepackage{array,booktabs,longtable}
\usepackage{xcolor}
\usepackage{tikz}
\usepackage[colorlinks=true,linkcolor=blue!55!black]{hyperref}
\setlength{\parindent}{0pt}
\setlength{\parskip}{7pt}
\renewcommand{\arraystretch}{1.27}
\newcommand{\err}{\textcolor{red}{\textbf{Ошибка:}}}
\newcommand{\idea}{\textcolor{blue!60!black}{\textbf{Ключевая идея:}}}
\newcommand{\outcome}{\textcolor{teal!60!black}{\textbf{Итог:}}}
\begin{document}
\begin{titlepage}
\centering
\vspace*{0.30\textheight}
{\Huge\bfseries Ревью домашней работы\par}
\vspace{2cm}
{\Large Автор: Лёвин Артём Александрович\par}
\vspace{0.6cm}
{\large 10 октября 2026 года\par}
\vfill
\begin{tikzpicture}
\draw[blue!35!black,line width=0.7pt] (-3,0) .. controls (-1.6,0.65) and (1.6,-0.65) .. (3,0);
\end{tikzpicture}
\vspace*{1cm}
\end{titlepage}

\section*{Сводная таблица заданий с замечаниями}
\small
\setlength{\tabcolsep}{3.5pt}
\begin{longtable}{>{\centering\arraybackslash}p{0.065\linewidth} >{\raggedright\arraybackslash}p{0.19\linewidth} >{\raggedright\arraybackslash}p{0.27\linewidth} >{\raggedright\arraybackslash}p{0.18\linewidth} >{\raggedright\arraybackslash}p{0.18\linewidth}}
\toprule
\textbf{№} & \textbf{Тема} & \textbf{Суть замечания} & \textbf{Ваш ответ} & \textbf{Правильный ответ}\\
\midrule
\endfirsthead
\toprule
\textbf{№} & \textbf{Тема} & \textbf{Суть замечания} & \textbf{Ваш ответ} & \textbf{Правильный ответ}\\
\midrule
\endhead
\hyperref[task:2]{2} & Вершина и сдвиг графика модуля & Не указана вершина и направление сдвига & «Сдвиг на 2 по оси~$x$» & $(-2;0)$, на 2 влево\\
\hyperref[task:7]{7} & Модуль квадратного трёхчлена & Неверный знак перед $x$ во второй ветви & $-x^2+x+2$ & $-x^2-x+2$\\
\hyperref[task:8]{8} & Произведение с $|x|$ & Появилась посторонняя переменная $y$; потеряна единица & $x(y-4)+1$; $-x(x-4)$ & $x^2-4x+1$; $-x^2+4x+1$\\
\hyperref[task:9]{9} & Произведение с $|x-4|$ & Ошибочный знак в первой промежуточной формуле; лишний модуль во второй & $x(x+4)+1$; затем правильное раскрытие & $x^2-4x+1$; $-x^2+4x+1$\\
\hyperref[task:10]{10} & График, вершина и нули & Неверные значения левой ветви; неправильный график, нули не записаны & Левая таблица: $(-1;-4)$, $(-2;-5)$ & Вершина $(-1;-2)$; нули $-3$ и $1$\\
\bottomrule
\end{longtable}
\setlength{\tabcolsep}{6pt}
\normalsize

\clearpage
\vspace*{\fill}
\begin{center}
{\Huge\bfseries Разбор ошибок}
\end{center}
\vspace*{\fill}
\clearpage

\phantomsection\label{task:2}
\section*{Задание № 2. Вершина и сдвиг графика}
\textbf{Текст задания.} Найдите вершину графика $y=|x+2|$. Укажите его сдвиг относительно графика $y=|x|$.

\textbf{Ваше решение.} Вы правильно раскрыли модуль:
\[
y=\begin{cases}x+2,&x\geqslant-2,\\-x-2,&x<-2.\end{cases}
\]
После этого записано: «сдвиг на 2 по оси $x$».

\textbf{Ваш ответ.} Сдвиг на две единицы по горизонтальной оси. Координаты вершины и направление сдвига отдельно не указаны.

\err\ Неполный ответ на два требования условия.

\textbf{Где именно возникает недочёт.} Последний правильный шаг --- разбиение функции на две ветви с границей $x=-2$. Далее следовало назвать точку встречи ветвей и уточнить, в какую сторону сместился график. Эти элементы отсутствуют.

\textbf{Почему этого недостаточно.} Вершина графика модуля --- точка, где выражение под знаком модуля равно нулю. Сдвиг «на две единицы по оси $x$» не определяет направление: возможен сдвиг как вправо, так и влево. В данном случае знак \emph{плюс} внутри модуля соответствует сдвигу \emph{влево}, а не вправо.

\idea\ Для функции $y=|x-a|+b$ вершина имеет координаты $(a;b)$. Здесь $x+2=x-(-2)$, поэтому $a=-2$, $b=0$.

\textbf{Исправление от последнего правильного шага.} По условию $x+2=0$ получаем $x=-2$. Подставляем: $y=|-2+2|=0$. Вершина --- $(-2;0)$. Чтобы получить $|x+2|$ из $|x|$, каждую точку исходного графика нужно перенести влево на две единицы.

\textbf{Полное правильное решение.} Начальный график $y=|x|$ имеет вершину $(0;0)$. Для нового графика имеем
\[
|x+2|=\begin{cases}x+2,&x\geqslant-2,\\-x-2,&x<-2.\end{cases}
\]
Правая ветвь начинается при $x=-2$ со значения $0$. Левая ветвь при приближении $x$ к $-2$ также подходит к значению $0$. Обе ветви встречаются в точке $(-2;0)$. Значит, график сохраняет форму буквы V, а его вершина переместилась из $(0;0)$ в $(-2;0)$.

\textbf{Проверка.} $y(-2)=0$, $y(-3)=1$, $y(-1)=1$. Точки $(-3;1)$ и $(-1;1)$ равноудалены от прямой $x=-2$, что подтверждает положение вершины.

\outcome\ Вершина $(-2;0)$; сдвиг относительно $y=|x|$ на две единицы \textbf{влево}.

\textbf{Задача на закрепление.} Найдите вершину графика $y=|x-5|$ и укажите его сдвиг относительно $y=|x|$.

\clearpage
\phantomsection\label{task:7}
\section*{Задание № 7. Модуль квадратного трёхчлена}
\textbf{Текст задания.} Представьте функцию $y=|x^2+x-2|$ без знака модуля с условиями для каждой ветви.

\textbf{Ваше решение.} Вначале правило модуля записано правильно:
\[
y=\begin{cases}x^2+x-2,&x^2+x-2\geqslant0,\\-(x^2+x-2),&x^2+x-2<0.\end{cases}
\]
Затем для второй ветви записано $y=-x^2+x+2$.

\textbf{Ваш ответ.} При $x^2+x<2$: $y=-x^2+x+2$.

\err\ При снятии скобок после минуса изменён знак не у всех слагаемых.

\textbf{Где именно возникает ошибка.} Последний правильный шаг: $-(x^2+x-2)$. Первый неверный шаг: $-x^2+x+2$. Перед слагаемым $x$ ошибочно остался знак плюс.

\textbf{Почему это неверно.} Минус перед скобками означает умножение \emph{каждого} слагаемого на $-1$. Поэтому $-(x^2+x-2)=-x^2-x+2$. Меняются три знака: перед $x^2$, перед $x$ и перед числом $-2$. Если хотя бы один знак не поменять, получится уже другая функция.

\idea\ Сначала обозначьте выражение под модулем как единое целое. Если оно неотрицательно, оставьте его без изменений; если отрицательно --- поставьте минус перед \emph{всем} выражением и лишь затем раскройте скобки.

\textbf{Исправление от последнего правильного шага.}
\[
-(x^2+x-2)=(-1)x^2+(-1)x+(-1)(-2)=-x^2-x+2.
\]

\textbf{Полное правильное решение.} Положим $A=x^2+x-2$. По определению модуля $|A|=A$ при $A\geqslant0$ и $|A|=-A$ при $A<0$. Тогда
\[
y=\begin{cases}
x^2+x-2,&x^2+x-2\geqslant0,\\
-x^2-x+2,&x^2+x-2<0.
\end{cases}
\]
Условия можно также записать как $x^2+x\geqslant2$ и $x^2+x<2$. Решать эти квадратные неравенства по условию задания не требуется.

\textbf{Проверка.} Возьмём $x=\tfrac12$. Подмодульное выражение $\tfrac14+\tfrac12-2=-\tfrac54$ отрицательно, поэтому используем вторую ветвь. По исходной функции $y=\bigl|-\tfrac54\bigr|=\tfrac54$. По исправленной ветви $y=-\tfrac14-\tfrac12+2=\tfrac54$. В ошибочной записи получилось бы $-\tfrac14+\tfrac12+2=\tfrac94$, что противоречит исходной функции.

\outcome\ При $x^2+x-2<0$ правильная ветвь --- $y=-x^2-x+2$.

\textbf{Задача на закрепление.} Представьте без знака модуля функцию $y=|x^2+3x-4|$, указав условие для каждой ветви. Квадратные неравенства решать не требуется.

\clearpage
\phantomsection\label{task:8}
\section*{Задание № 8. Произведение с модулем}
\textbf{Текст задания.} Представьте функцию $y=|x|(x-4)+1$ без знака модуля с условиями для каждой ветви.

\textbf{Ваше решение.} Первое раскрытие верно:
\[
y=\begin{cases}x(x-4)+1,&x\geqslant0,\\-x(x-4)+1,&x<0.\end{cases}
\]
Однако ниже для первой ветви записано $y=x(y-4)+1$, а для второй --- равенство $-x(x-4)=y$ без прибавления единицы.

\textbf{Ваш ответ.} $x(y-4)+1$ при $x\geqslant0$; $y=-x(x-4)$ при $x<0$.

\err\ При переписывании правильных ветвей переменная $x$ заменена на $y$ внутри скобок и потеряно слагаемое $+1$.

\textbf{Где именно возникают ошибки.} Последний полностью верный шаг --- первая запись кусочной функции. Первый неверный переход --- $x(x-4)+1\longrightarrow x(y-4)+1$. Отдельно во второй ветви переход $-x(x-4)+1\longrightarrow -x(x-4)$ удаляет единицу.

\textbf{Почему это неверно.} В исходной формуле множитель $(x-4)$ зависит от аргумента $x$. Замена на $(y-4)$ делает равенство уравнением с $y$ в обеих частях, а не заданной функцией от $x$. Слагаемое $+1$ находится вне произведения, поэтому раскрытие модуля его не уничтожает; оно должно присутствовать в обеих ветвях.

\idea\ Раскрытие $|x|$ изменяет только множитель $|x|$: при $x\geqslant0$ он равен $x$, при $x<0$ равен $-x$. Остальные множители и прибавления сохраняются.

\textbf{Исправление от последнего правильного шага.} Перепишите две уже полученные верные строки и раскройте скобки:
\[
\begin{aligned}
x(x-4)+1&=x^2-4x+1,\\
-x(x-4)+1&=-x^2+4x+1.
\end{aligned}
\]

\textbf{Полное правильное решение.} При $x\geqslant0$ имеем $|x|=x$, значит $y=x^2-4x+1$. При $x<0$ имеем $|x|=-x$, значит $y=-x^2+4x+1$. Таким образом,
\[
\boxed{y=\begin{cases}
x^2-4x+1,&x\geqslant0,\\
-x^2+4x+1,&x<0.
\end{cases}}
\]
Граница между ветвями --- число $0$; оно отнесено к первой ветви, чтобы каждое значение аргумента имело ровно одно правило вычисления.

\textbf{Проверка.} При $x=2$ исходная формула даёт $2(2-4)+1=-3$; первая ветвь даёт $4-8+1=-3$. При $x=-1$ исходная формула даёт $1(-1-4)+1=-4$; вторая ветвь даёт $-1-4+1=-4$. Если забыть $+1$, вышло бы $-5$, то есть неверное значение. При $x=0$ получаем $y=1$.

\outcome\ Обе ветви содержат слагаемое $+1$; в формулах справа остаётся только переменная $x$.

\textbf{Задача на закрепление.} Представьте без знака модуля функцию $y=|x|(x+3)-2$ с условиями для каждой ветви.

\clearpage
\phantomsection\label{task:9}
\section*{Задание № 9. Произведение с $|x-4|$}
\textbf{Текст задания.} Представьте функцию $y=x|x-4|+1$ без знака модуля с условиями для каждой ветви.

\textbf{Ваше решение.} Граница $x=4$ определена правильно. Однако в первой кусочной записи при $x\geqslant4$ стоит $x(x+4)+1$, а во второй ветви сохраняется модуль в записи $x(-|x-4|)+1$. Ниже Вы уже правильно записали итоговые многочлены $x^2-4x+1$ и $-x^2+4x+1$.

\textbf{Ваш ответ.} В нижней записи указаны правильные ветви: $x^2-4x+1$ и $-x^2+4x+1$. Однако промежуточные равенства им противоречат.

\err\ Неверное раскрытие модуля в промежуточной записи; последующие правильные строки не вытекают из написанного шага.

\textbf{Где именно возникает ошибка.} Последний правильный шаг --- условие $x\geqslant4$ для неотрицательного $x-4$. Первый неправильный шаг --- замена $|x-4|$ на $x+4$. Также при $x<4$ нужно сразу заменить $|x-4|$ на $-(x-4)$, а не писать $-|x-4|$.

\textbf{Почему это неверно.} Знак внутри модуля не меняется произвольно. Если $x-4\geqslant0$, то $|x-4|=x-4$, со знаком \emph{минус} перед четырьмя. Если $x-4<0$, то $|x-4|=-(x-4)=4-x$. Сохранение модуля с добавленным внешним минусом тоже не является его раскрытием: например, при $x=3$ число $-|3-4|=-1$, тогда как $|3-4|=1$.

\idea\ В произведении $x|x-4|$ раскрывайте только множитель $|x-4|$, а множитель $x$ оставляйте без изменений. Условие ветви определяется знаком $x-4$, а не знаком $x$.

\textbf{Исправление от последнего правильного шага.}
\[
\begin{aligned}
x\geqslant4:&\quad x|x-4|+1=x(x-4)+1=x^2-4x+1,\\
x<4:&\quad x|x-4|+1=x(4-x)+1=-x^2+4x+1.
\end{aligned}
\]

\textbf{Полное правильное решение.} Сначала находим, где подмодульное выражение равно нулю: $x-4=0$, значит $x=4$. Справа от этой точки выражение $x-4$ неотрицательно; слева оно отрицательно. Поэтому
\[
\boxed{y=\begin{cases}
x^2-4x+1,&x\geqslant4,\\
-x^2+4x+1,&x<4.
\end{cases}}
\]
Это совпадает с Вашей нижней записью. Исправлять следует прежде всего обоснование перехода к ней: правильный ответ должен следовать из правильных равенств.

\textbf{Проверка.} При $x=4$ исходное значение равно $4\cdot0+1=1$, и первая ветвь тоже даёт $16-16+1=1$. Ошибочная промежуточная формула дала бы $4(4+4)+1=33$. При $x=3$ исходная функция даёт $3\cdot1+1=4$; вторая исправленная ветвь --- $-9+12+1=4$.

\outcome\ Итоговые многочлены были записаны верно; ошибочные промежуточные преобразования необходимо заменить корректным раскрытием модуля.

\textbf{Задача на закрепление.} Представьте без знака модуля функцию $y=x|x+2|-3$ с условиями для каждой ветви.

\clearpage
\phantomsection\label{task:10}
\section*{Задание № 10. Кусочная функция и её график}
\textbf{Текст задания.} Для $y=|x+1|-2$ запишите кусочную формулу, найдите вершину и нули. Постройте график и проверьте значение в точке смены ветвей.

\textbf{Ваше решение.} Формулы ветвей получены правильно:
\[
y=\begin{cases}x-1,&x\geqslant-1,\\-x-3,&x<-1.\end{cases}
\]
Для первой ветви составлена верная таблица $(-1;-2)$, $(1;0)$. Во второй таблице стоят значения $(-1;-4)$ и $(-2;-5)$; далее левая часть графика проведена как возрастающая прямая. Нули и координаты вершины отдельным ответом не записаны.

\textbf{Ваш ответ.} Правильная кусочная формула, но неверные точки и линия левой ветви; отсутствуют отдельные ответы о вершине, нулях и проверке в точке соединения.

\err\ Ошибки при подстановке отрицательных значений аргумента во вторую ветвь и при построении графика; не выполнены все пункты условия.

\textbf{Где именно возникает ошибка.} Последний правильный шаг --- запись $y=-x-3$ при $x<-1$. Первый неверный численный шаг --- строка второй таблицы, где при $x=-2$ указано $y=-5$. Верно: $-(-2)-3=-1$. Кроме того, $x=-1$ вообще не входит в условие $x<-1$; в точке границы следует пользоваться первой ветвью.

\textbf{Почему это неверно.} Для отрицательного аргумента $x=-2$ выражение $-x$ положительно и равно $2$. Поэтому $-x-3$ при движении по оси $x$ вправо \emph{убывает}: коэффициент при $x$ равен $-1$. Изображённая левая линия возрастает, то есть имеет противоположный наклон. Две ветви графика модуля должны встретиться в одной вершине без разрыва.

\idea\ При построении кусочной функции проверяйте сразу три вещи: принадлежность точки условию ветви, правильность арифметики и направление наклона прямой. Для $|x-a|+b$ обе части графика сходятся в точке $(a;b)$.

\textbf{Исправление от последнего правильного шага.} Сохраняем найденные формулы. Для левой ветви при $x=-3$ имеем $y=3-3=0$, а при $x=-2$ получаем $y=2-3=-1$. На границе $x=-1$ получаем $y=-2$. Поэтому левая линия должна проходить через точки $(-3;0)$, $(-2;-1)$ и подходить к $(-1;-2)$.

\textbf{Полное правильное решение.} Раскрываем модуль по знаку $x+1$:
\[
y=\begin{cases}
x+1-2=x-1,&x\geqslant-1,\\
-(x+1)-2=-x-3,&x<-1.
\end{cases}
\]
Подмодульное выражение обращается в ноль при $x=-1$. Подставляя это значение, находим $y=|-1+1|-2=-2$. Следовательно, \textbf{вершина} --- $(-1;-2)$.

\textbf{Нули функции.} Приравниваем выражение к нулю:
\[
|x+1|-2=0\quad\Longleftrightarrow\quad |x+1|=2.
\]
Значит, $x+1=2$ или $x+1=-2$, откуда $x=1$ или $x=-3$. Оба значения подходят. Для правой ветви $1-1=0$; для левой $-(-3)-3=0$.

\textbf{Опорные точки для графика.} Левая ветвь: $(-5;2)$, $(-3;0)$, $(-2;-1)$. Точка $(-1;-2)$ является общей вершиной и по правилу кусочной функции относится к правой ветви. Правая ветвь: $(-1;-2)$, $(1;0)$, $(3;2)$. Соединяя точки прямыми лучами, получаем V-образный график.

\textbf{Проверка в точке смены ветвей.} При $x=-1$ исходная функция даёт $|-1+1|-2=-2$. Правая формула даёт $-1-1=-2$. Если подставить $x=-1$ в продолжение левой формулы, получим $-(-1)-3=-2$; значения сходятся. При этом сама левая ветвь определена строго для $x<-1$.

\outcome\ Вершина $(-1;-2)$; нули $x=-3$ и $x=1$; левая ветвь убывает при движении вправо, правая возрастает. График непрерывен в точке $x=-1$.

\textbf{Задача на закрепление.} Для функции $y=|x-2|-3$ запишите кусочную формулу, найдите вершину и нули, постройте график и проверьте значение в точке смены ветвей.

\clearpage
\begin{center}
\large\bfseries Исправленный график к заданию № 10
\end{center}
\vspace{1cm}
\begin{center}
\begin{tikzpicture}[scale=0.96]
\draw[step=1cm,gray!17,very thin] (-5.7,-3.25) grid (3.7,2.7);
\draw[->] (-5.8,0) -- (3.85,0) node[right] {$x$};
\draw[->] (0,-3.3) -- (0,2.95) node[above] {$y$};
\foreach \x in {-5,-4,-3,-2,-1,1,2,3} {\draw (\x,0.07)--(\x,-0.07) node[below] {\scriptsize $\x$};}
\foreach \y in {-3,-2,-1,1,2} {\draw (0.07,\y)--(-0.07,\y) node[left] {\scriptsize $\y$};}
\draw[very thick,blue!65!black] (-5.5,2.5) -- (-1,-2) -- (3.5,2.5);
\fill (-3,0) circle(2.5pt) node[above left] {\small $(-3;0)$};
\fill (-1,-2) circle(2.5pt) node[below right] {\small $(-1;-2)$};
\fill (1,0) circle(2.5pt) node[below right] {\small $(1;0)$};
\end{tikzpicture}
\end{center}
\vfill
\begin{center}
Лёвин Артём Александрович эксклюзивно для Марины Селиверстовой
\end{center}
\end{document}
